\chapter{Bursitis}

\section*{Definition}
Bursitis is inflammation of a bursa, a fluid-filled sac that reduces friction between tendons, muscles, and bone. It presents with localized pain, swelling, and tenderness over affected sites such as the shoulder (subacromial), elbow (olecranon), hip (trochanteric), knee (prepatellar), or heel (retrocalcaneal).

\section*{What Happens in Your Body}
Repetitive microtrauma or acute injury irritates the synovial lining of the bursa, triggering an inflammatory cascade involving prostaglandins, leukotrienes, and immune signaling proteins. This results in increased synovial fluid production, bursal distension, and pain with movement. Chronic bursitis may lead to fibrous thickening and adhesion formation.

\section*{Causes}
\begin{itemize}
 \item \textbf{Mechanical:} Repetitive friction or pressure (e.g., kneeling, leaning on elbows, throwing sports)
 \item \textbf{Traumatic:} Direct blow or fall onto the bursa
 \item \textbf{Infectious:} Septic bursitis from bacterial inoculation (often \textit{S. aureus}) through skin breach
 \item \textbf{Inflammatory:} Gout, pseudogout, rheumatoid arthritis, systemic lupus erythematosus
 \item \textbf{Occupational:} Prolonged kneeling (housemaid's knee), carpet laying, gardening
\end{itemize}

\section*{When to See a Doctor}

\begin{itemize}
 \item Overlying erythema, warmth, or fever suggesting septic bursitis
 \item Pain severe enough to limit range of motion or daily activities
 \item No improvement after 5--7 days of conservative management
 \item Recurrent bursitis in the same location
\end{itemize}

\section*{Related Symptoms}
\begin{itemize}
 \item Localized swelling and fluctuance over the bursa
 \item Tenderness with direct palpation
 \item Pain with motion of the adjacent joint
 \item Erythema and warmth (suggests infection)
\end{itemize}
\textbf{Important:} Septic bursitis requires urgent aspiration and antibiotics; do not inject corticosteroids until infection is excluded.

\section*{Self-Care / Home Management}
\begin{itemize}
 \item Rest and avoidance of aggravating activities
 \item Ice packs applied for 15--20 minutes every 2--3 hours during the acute phase
 \item Compression with an elastic bandage to limit swelling
 \item Elevation of the affected limb when possible
 \item Over-the-counter NSAIDs (ibuprofen, naproxen) for pain and inflammation
\end{itemize}

\section*{How Your Doctor Will Evaluate You}
\begin{itemize}
 \item Clinical diagnosis based on history and physical examination
 \item Ultrasound to confirm bursal fluid and guide aspiration
 \item Bursal fluid analysis: cell count, Gram stain, culture, crystal analysis
 \item MRI if deeper bursae (e.g., trochanteric) are suspected or diagnosis is unclear
 \item Radiography to exclude fracture or joint pathology
\end{itemize}

\section*{Treatment Options}
\begin{itemize}
 \item Activity modification and physical therapy for biomechanical correction
 \item NSAIDs for 7--14 days for pain and inflammation
 \item Intrabursal corticosteroid injection for non-septic bursitis
 \item Antibiotics (empiric coverage for \textit{S. aureus} and \textit{Streptococcus}) for septic bursitis
 \item Bursectomy for refractory cases after failure of conservative measures
\end{itemize}

\section*{References}
\begin{thebibliography}{99}
\bibitem{bursitis:ref1} Aaron DL, Patel A, Kayiaros S, et al. ``Four Common Types of Bursitis: Diagnosis and Management.'' J Am Acad Orthop Surg. 2011;19(6):359--367.
\bibitem{bursitis:ref2} Darnell M, Ho SS, Malanga GA. ``Ultrasound-Guided Bursal Injection.'' Phys Med Rehabil Clin N Am. 2016;27(3):685--695.
\bibitem{bursitis:ref3} Baumbach SF, Krapohl BD, Hungerer SA, et al. ``Prepatellar and Olecranon Bursitis: Literature Review and Development of a Treatment Algorithm.'' Arch Orthop Trauma Surg. 2014;134(3):359--370.
\bibitem{bursitis:ref4} Reilly JP, Lally JF, Nolla JM, et al. ``Septic Bursitis: Diagnosis and Treatment.'' Rheum Dis Clin North Am. 2020;46(2):287--300.
\bibitem{bursitis:ref5} Butcher JD, Salzman KL, Lillegard WA. ``Trochanteric Bursitis: A Review.'' Am Fam Physician. 1996;54(4):1243--1250.
\end{thebibliography}
